\chapter{Understanding the Network-Level Behavior of Spammers}
\paperinfo{A. Ramachandran, N. Feamster (Georgia Tech)}{ACM SIGCOMM 2006}

\begin{keybox}[The paper in three sentences]
Spam filters mostly look at the \emph{content} of an email, which spammers change freely. This paper looks instead at
\key{network-level properties} of the sender (its IP address, IP range, AS, operating system, BGP route), which are much harder to
fake, by joining 17 months of spam caught at a \key{sinkhole} domain with BGP data and a botnet trace. It finds that spam comes from a
few concentrated IP ranges and ASes, mostly from \key{short-lived Windows bots} that send very little each, and that a small group of
spammers \key{hijacks large IP blocks with short-lived BGP announcements} (``\key{BGP spectrum agility}'') to stay untraceable.
\end{keybox}

% ------------------------------------------------------------------
\section{The problem: filters look at the wrong thing}

There are two common ways to stop spam:
\begin{itemize}
  \item \key{Content filtering} (e.g.\ SpamAssassin): look at words, headers and patterns in the email. Spammers simply change the text,
        add random characters, use images, and so on. Changing content costs the spammer nothing, so this is an endless arms race.
  \item \key{IP blacklists (DNSBLs)}: lists of IP addresses known to send spam. The mail server looks up the connecting IP and rejects it if
        it is listed. This only works if a spammer keeps using the same IP for a while.
\end{itemize}

The paper's argument: a spammer can rewrite an email in a second, but it is much harder to change \emph{where the email is sent from}: the
ISP, the IP address block, the AS, the botnet. So if spam differs from normal mail in these \key{network-level} properties, filters built on
them would be more robust. Before building such filters, we need to know how spammers actually behave at the network level. That is what
this paper measures.

\textbf{Which properties can be trusted?} Most email headers can be forged. But the \key{IP address of the machine that opens the TCP
connection} to our mail server cannot be forged (the TCP handshake needs replies to reach it), and from it we can derive its AS, its
route, and a TCP fingerprint. The whole study uses only such unforgeable properties.

% ------------------------------------------------------------------
\section{How spammers send spam (the methods being measured)}

\begin{center}
\small
\begin{tabularx}{\linewidth}{@{}lX@{}}
\toprule
Method & How it works\\
\midrule
\key{Direct spamming} & Buy connectivity from a ``spam-friendly'' ISP and send directly. If the ISP kicks them out, they must move and renumber.\\
\key{Open relays / proxies} & Mail servers or proxies that forward mail for anyone. They hide the spammer behind one extra hop. Blacklists have
nearly killed this method.\\
\key{Botnets} & Thousands of infected home PCs (\emph{drones}) controlled by one master. Each drone sends spam from its own IP. The paper studies
\key{Bobax}, a Windows worm (spread through DCOM/LSASS holes) that turns a PC into a mail relay.\\
\key{BGP spectrum agility} (new) & Announce a route for an IP block you do not own (often a huge /8), send spam from addresses in it, then
withdraw the route. Described below.\\
\bottomrule
\end{tabularx}
\end{center}

% ------------------------------------------------------------------
\section{The measurement setup: four datasets joined together}

\subsection{1. The spam sinkhole}
The authors registered a domain that has \emph{no real users}, and published an MX record for it. Any mail sent there must be spam, because
no-one legitimate has an address there. Spammers find the domain through address harvesting and dictionary attacks. The sinkhole collected
everything from August 2004 to January 2006: over \key{10 million spam messages}.

The mail server runs \key{Mail Avenger}, a programmable SMTP server. It accepts every message (whatever the username), and \emph{at the moment
the sender connects} it records:
\begin{itemize}
  \item the \textbf{IP address} of the relay that made the SMTP connection;
  \item a \textbf{traceroute} to that IP (where it is in the network);
  \item a passive \key{p0f TCP fingerprint}: small details of the TCP SYN (initial window, TTL, options) that reveal the sender's operating
        system, without sending any probes;
  \item \textbf{DNSBL lookups} in 8 blacklists, \emph{at that moment} (not later, when the IP might have been added).
\end{itemize}

\subsection{2. Legitimate mail logs}
To see whether spam looks different from normal mail, they used one day of logs (about 700,000 accepted messages) from a large email provider
with about 40 million mailboxes. The provider rejects likely spam, so the accepted mail is a sample of legitimate senders.

\subsection{3. A botnet's membership list (Bobax trace)}
Bobax drones contacted their controller by a \emph{domain name}. Researchers got the domain's DNS operators to point that name at a machine they
controlled (a \key{C\&C sinkhole}). For 1.5 days every drone that phoned home connected to the researchers instead, so they captured the IP
addresses of \key{117,268 infected hosts}. Joining this list with the sinkhole tells us which spam came from \emph{known} bots.

\subsection{4. BGP routing data}
A \key{BGP monitor} in the same network as the sinkhole had an iBGP session with the border router. It saw every change in the border router's
best route to each prefix. So for every spam message, the authors could ask: \emph{at the time this spam arrived, was there a route to the
sender's IP, and how long did that route exist?}

\textbf{Limitation (the authors stress it)}: this is one domain's view, so the numbers are not Internet-wide. But it is exactly the view a
single network operator would have, which is what a real filter would see.

% ------------------------------------------------------------------
\section{Finding 1: where spam comes from in IP space and ASes}

\paperfig[0.48]{sp2_fig2}{CDF over IP space (binned by /24) of legitimate email, spam, and spamming IP addresses. Flat parts are ranges that send
nothing; steep parts are ranges that send a lot.\label{fig:sp2ip}}

\textbf{How to read Figure~\ref{fig:sp2ip}.} The x-axis is the whole IPv4 space from 0.0.0.0 to 255.255.255.255. The curve rises wherever mail
comes from. If spam were spread evenly over the address space, the curve would be a straight diagonal line. Instead it has steep jumps and long flat
sections.

\begin{itemize}
  \item \textbf{Finding 4.1:} most spam comes from a \key{small part of IP space} (for example 60.*--90.* and 200.*--220.*); large ranges send
        none at all.
  \item But \textbf{legitimate mail comes from mostly the same ranges}. A few ranges differ: 80.*--90.* sends mostly spam, while 60.*--70.* sends
        relatively more legitimate mail. So IP range alone separates spam from good mail only for some ranges.
  \item \textbf{Ranges are persistent; individual IPs are not.} The per-range distribution looks the same month after month. But about
        \key{85\% of spamming IPs sent fewer than 10 emails} to the sinkhole in 17 months. Blocking one IP does little; noticing that a whole
        range or AS keeps spamming is more useful.
  \item \textbf{Finding 4.2 (ASes):} two ASes (Korean Internet Exchange and China Telecom) sent about 12\% of all spam; the top 20 ASes sent about
        37\%. These ASes are mostly \emph{different} from the top ASes for legitimate mail (Google, Yahoo, AT\&T, \dots). So ``which AS sent
        this'' is a useful hint.
  \item \textbf{Finding 4.3 (countries):} the top two ASes are in Asia, but 11 of the top 20 are in the US, and overall about three times as much
        spam comes from US ISPs as from Korea or China.
\end{itemize}

\subsection{How well do blacklists work?}
Because Mail Avenger queried 8 DNSBLs at the time each spam arrived, the authors can measure this directly.
\begin{itemize}
  \item About \key{80\%} of spam came from an IP listed in \emph{at least one} of the 8 lists. That sounds good, but:
  \item the most aggressive single list (Spamcop) listed only about \key{half} of the spam, and aggressive lists also have many false positives;
  \item spam from \key{short-lived BGP routes} was listed much less: \key{half} of those IPs were in \emph{no} list at all.
\end{itemize}
Bobax drones were listed \emph{more} often than average, which seems odd for IPs that each sent so little to this domain. The reason: the same drone
spams many other domains too, and those reports reach the blacklists. This is evidence that \emph{collaborative} filtering across domains works.

% ------------------------------------------------------------------
\section{Finding 2: spam from botnets}

\subsection{Bots and spam live in the same IP ranges}
The IP distribution of Bobax drones looks almost the same as the distribution of all spamming hosts in Figure~\ref{fig:sp2ip}. This is correlation,
not proof, but it suggests that much of the spam comes from botnets. One exception: the range 60.*--67.* has many spamming IPs but few Bobax drones,
which hinted at a different method there (spectrum agility, below).

\subsection{Operating systems (from p0f fingerprints)}
\begin{itemize}
  \item About \key{70\%} of spamming hosts (95\% of those that could be identified) run \key{Windows}, consistent with infected home PCs.
  \item \textbf{Finding 5.2:} only about 4\% of hosts run something else (mostly Linux/FreeBSD), but they send about \key{8\%} of the spam: each
        non-Windows host sends more. These are likely dedicated spam servers.
\end{itemize}

\subsection{How long does a bot spam, and how much?}
Only 4,693 of the 117,268 known Bobax drones ever spammed this one domain, but that sample tells us how bots behave.

\paperfig[0.48]{sp2_fig7}{Lifetime of each Bobax drone at the sinkhole (time between its first and last spam). About 65\% have lifetime 0: they
appeared once.\label{fig:sp2life}}

\begin{itemize}
  \item \textbf{Finding 5.3 (single-shot bots):} more than \key{65\%} of Bobax IPs spammed the domain \emph{only once}, and about \key{75\%} were active
        for \key{less than two minutes} in total. The rest are bimodal: they last a day or longer, some for months.
  \item \textbf{Finding 5.4 (low volume):} \key{99\%} of bots sent \key{fewer than 100} spams to this domain over 17 months, whether short- or long-lived.
        That is less than one spam per bot per day.
\end{itemize}

\textbf{Why this matters for filters.} A spammer with 100,000 bots can send a few messages from each and never reuse an address on the same
domain. Then:
\begin{itemize}
  \item \textbf{per-IP blacklisting} at one domain fails (the IP is never seen again);
  \item \textbf{volume-based detection} fails (each bot is quiet).
\end{itemize}
Only aggregate views (many domains sharing data, or looking at ranges/ASes/botnet membership) can catch this.

% ------------------------------------------------------------------
\section{Finding 3: BGP spectrum agility}

\subsection{What the attack is}
BGP has no check that an AS is allowed to announce a prefix. A spammer who controls a BGP-speaking router (for example at a small, careless ISP)
can do the following:
\begin{enumerate}
  \item \textbf{Announce} a route for an IP block it does not own, usually a large one such as 61.0.0.0/8, often \emph{unused} (``dark'') space.
        Now the Internet sends traffic for those addresses to the spammer.
  \item \textbf{Send spam} from many different IP addresses scattered across that block. The TCP handshake works, because replies to those
        addresses now reach the spammer.
  \item \textbf{Withdraw} the route a few minutes later. The addresses become unreachable again.
\end{enumerate}

\paperfig[0.5]{sp2_fig9}{One real episode for 61.0.0.0/8 on 30 Sep 2005: the route is announced at about 13:23, spam arrives from addresses
in the block over the next minutes, and the route is withdrawn at about 13:35.\label{fig:sp2agile}}

\textbf{Walking through Figure~\ref{fig:sp2agile}}:
\begin{center}
\small
\begin{tabular}{@{}ll@{}}
\toprule
Time (30 Sep 2005) & Event\\
\midrule
13:23 & AS 4678 announces 61.0.0.0/8 (a /8 is 16.7 million addresses)\\
13:24--13:32 & spam arrives at the sinkhole from relays with addresses inside 61.0.0.0/8\\
13:35 & the route is withdrawn; those addresses are unreachable again\\
later & a user complains; the source address cannot be reached or traced\\
\bottomrule
\end{tabular}
\end{center}

\subsection{Why it is so effective (for the spammer)}
\begin{itemize}
  \item \textbf{Untraceable}: by the time someone investigates, the route is gone, so the sender's network does not exist any more. Only someone who
        logged BGP updates at that exact time could connect the two.
  \item \textbf{Defeats blacklists}: a /8 gives millions of fresh addresses, so the spammer can ``hop'' to a new address for every message.
        Half of these IPs were in no blacklist.
  \item \textbf{Large prefixes get through}: ISPs commonly filter announcements of very small, specific prefixes, but rarely filter a /8.
  \item \textbf{Hides the origin}: many AS paths in these announcements contained reserved (unallocated) AS numbers, suggesting forged paths.
\end{itemize}

\textbf{The most persistent cases}: AS 21562 (an ISP in Indianapolis) announcing 66.0.0.0/8; AS 8717 (Sofia, Bulgaria) announcing 82.0.0.0/8; and
AS 4678 (Canon Network Communications, Japan) announcing 61.0.0.0/8. In February 2006, 60--80\% of the relays seen using this technique did not
respond to traceroute.

\subsection{How common is it?}
\begin{itemize}
  \item About 99\% of spam came from routes that had existed for at least a day, so spectrum agility is \emph{not} the main method.
  \item \textbf{Finding 6.2:} only about \key{1\%} of spam came from routes that lived less than a day; but in periods when the technique was in use,
        it was up to \key{10\%}.
  \item \textbf{Finding 6.3:} the ASes that use it \emph{most persistently} are not the ones that send the \emph{most} spam through it. Persistent users
        seem to do it deliberately; for high-volume ones the short route may just coincide with spam sent another way.
  \item It can be combined with other methods (for example, to give a botnet fresh addresses), which makes it more worrying than its volume suggests.
\end{itemize}

% ------------------------------------------------------------------
\section{Lessons for building better spam defences}

Network-level features have two advantages over content: they are \key{harder to change}, and they can be seen \key{inside the network}, close to the
source, before the mail reaches the destination server.

\begin{center}
\small
\begin{tabularx}{\linewidth}{@{}lX@{}}
\toprule
Lesson & Why (from which finding)\\
\midrule
1. Filtering needs a \key{better notion of host identity} than an IP address. & Single-shot bots, hijacked /8s and dynamic addressing all break
``IP = sender''.\\
2. \key{Aggregate behaviour} (ranges, ASes, botnet membership, across many domains) exposes spammers better than single IPs. & Individual IPs are
transient, but ranges and ASes are persistent; blacklists work for bots because many domains report them.\\
3. \key{Securing BGP} is needed for traceability. & Spectrum agility works only because anyone can announce any prefix.\\
4. Some network features are \key{easy to add} to filters. & E.g.\ ``was this sender's prefix announced only minutes ago?'' would stop spectrum
agility, which blacklists miss.\\
\bottomrule
\end{tabularx}
\end{center}

\begin{keybox}[Exam one-liner]
Using a 17-month spam sinkhole (a domain with no users, so all mail is spam) joined with p0f fingerprints, real-time DNSBL lookups, a Bobax botnet
membership trace and BGP updates, the paper shows: spam comes from a few persistent IP ranges and ASes, but individual IPs are transient; most spam comes
from Windows bots that appear once, for under 2 minutes, and send under 100 messages; and some spammers use \emph{BGP spectrum agility}: briefly announce
a hijacked /8, spam from scattered addresses in it, withdraw it, which defeats blacklists and tracing. So filters should use network-level, aggregate
features, and BGP must be secured.
\end{keybox}

\section{Revision checklist}
\begin{itemize}
  \item Why network-level: content is cheap to change; sender IP (TCP handshake), AS, route and OS fingerprint are hard to fake.
  \item Spamming methods: direct, open relays/proxies, botnets (Bobax), BGP spectrum agility.
  \item Data: sinkhole + Mail Avenger (IP, traceroute, p0f, 8 DNSBLs at receive time); legitimate mail logs; Bobax C\&C sinkhole (117K drones via DNS redirection); co-located BGP monitor.
  \item IP space: concentrated, similar to legitimate mail except some ranges (80--90.* mostly spam); ranges persistent, 85\% of IPs send $<$10 mails.
  \item ASes: top 2 ASes about 12\%, top 20 about 37\%; different from legitimate-mail ASes.
  \item Blacklists: 80\% listed in at least one of 8, best single list about 50\%, spectrum-agility IPs half unlisted.
  \item Bots: 70\% Windows; non-Windows 4\% of hosts send 8\% of spam; 65\% single-shot, 75\% active $<$2 min, 99\% send $<$100 spams.
  \item Spectrum agility: announce /8 $\to$ spam from scattered addresses $\to$ withdraw; untraceable, beats blacklists, /8s pass route filters; about 1\% of spam (up to 10\%).
  \item Lessons: better host identity, aggregate behaviour, secure BGP, add features like ``recently announced prefix''.
\end{itemize}
